\section{Common assumptions}
\label{sec:assumptions}
The common assumptions for both FE-MSM and SeqGPLVM-MSM are stated below. 
The assumptions can be categorized into two groups: identification assumptions and 
modeling assumptions. Identification assumptions, 
are the assumptions which are needed to write target causal estimands such as 
$\mathbb{E}[Y_i(\underline{\mathbf{a}}_{T-k})]$ in terms of observed data. 
Based on the identification strategy and the estimator of choice, 
certain quantities must be modeled. 
For example, estimating a marginal structural model for the target causal estimand using 
inverse probability of treatment weighting (IPTW) requires modeling the 
propensity scores. Modeling assumptions are about the 
requirements placed on these models in order to have unbiased or consistent estimates. 
Throughout the text, 
we use C, I, and M to denote common, identification, and modeling 
assumptions, respectively.
% --- Common assumptions (C1, C2, ...) ---
% Assumes you already defined the assumpC environment (C-prefix) in the preamble.

\begin{idassumpC}[Consistency]\label{ass:C-consistency}
The observed outcome for each individual under their observed treatment history equals
the potential outcome under that treatment history. Formally, if $\bar{A}_{iT}=\bar{\mathbf{a}}_{iT}$, 
then $Y_{i}=Y_{i}(\bar{\mathbf{a}}_{iT})$. 
Moreover, $Y_{i} = Y_{i}(\underline{\mathbf{a}}_{i,T-k})$ if 
$\bar{\mathbf{A}}_{iT}=\left(\bar{\mathbf{A}}_{i,T-k-1},\underline{\mathbf{a}}_{i,T-k}\right)$.
\end{idassumpC}

\begin{comment}

\begin{assumpC}[Positivity]\label{ass:C-positivity}
\[
\Pr\!\left[ A_t = a_t \,\middle|\, 
\bar{A}_{t-1} = \bar{a}_{t-1},\,
\bar{X}_t = \bar{x}_t \right] > 0
\]

for all values $(\bar{a}_{t-1}, \bar{x}_t)$ with

\[
\Pr\!\left[ \bar{A}_{t-1} = \bar{a}_{t-1},\,
\bar{X}_t = \bar{x}_t \right] \neq 0
\]

in the population of interest, and for all $a_t \in \mathcal{A}_t$.
\end{assumpC}

\end{comment}

\begin{idassumpC}[No interference]\label{ass:C-no-interference}
The treatment assignment for one individual does not affect the potential outcomes of
other individuals. This assumption is also known as the no spillover assumption.
\end{idassumpC}

\begin{idassumpC}[Time-invariance of unmeasured confounding]\label{ass:C-timeinvariant-uc}
The unmeasured confounding is assumed to be
time-invariant. That is, $\mathbf{U}_{i,t} = \mathbf{U}_i$ for $t \in \{1,\dots,T\}$. We can then write 
conditional ignorability as 
\[
\{ Y_{i,t+1}(\bar{\mathbf{a}}_{t}),\, \bar{\mathbf{X}}_{i,t+1}(\bar{\mathbf{a}}_{t}) \}
\perp\!\!\!\perp
A_{it}
\mid
\bar{\mathbf{X}}_{it},\,
\bar{\mathbf{A}}_{i,t-1} = \bar{\mathbf{a}}_{t-1},\,
\mathbf{U}_i .
\]
Note that this assumption is about the true underlying data-generating process and cannot be  
used for identification directly since $\mathbf{U}_i$ is unmeasured. 
\end{idassumpC}

\begin{idassumpC}[Positivity]\label{ass:C-positivity}
For all treatment times $t = T - k + 1, \ldots, T$, 
and all histories $(\bar{\mathbf{X}}_t, \bar{\mathbf{A}}_{t-1}, \boldsymbol{\eta}_i)$ 
with positive probability, every treatment level has positive probability:
\[
\mathbb{P}(A_t = a_t \mid \bar{\mathbf{X}}_t = \bar{\mathbf{x}}_t,\ \bar{\mathbf{A}}_{t-1} = 
\bar{\mathbf{a}}_{t-1}, \boldsymbol{\eta}_i) > 0 \quad \text{for all } a_t \in \mathcal{A}_t
\]
Where $\boldsymbol{\eta}_i$ represents the unit-specific fixed effect in the FE-MSM model 
or the confounder substitute in the SeqGPLVM model. This assumption  particularly
is necessary for the identification of $\mathbb{E}[Y_i(\underline{\mathbf{a}}_{T-k})]$.
\end{idassumpC}


